\documentclass[10pt,a4paper]{amsart}
\usepackage{amsmath,amssymb,amsthm,graphicx}
\usepackage{verbatim}
\newenvironment{reference}{\comment}{\endcomment}
%\newenvironment{reference}{}{}
\newenvironment{slogan}{\comment}{\endcomment}
\newenvironment{history}{\comment}{\endcomment}

% For commutative diagrams we use Xy-pic
\usepackage[all]{xy}

% We use 2cell for 2-commutative diagrams.
\xyoption{2cell}
% Theorem environments.
%
\theoremstyle{plain}
\newtheorem{theorem}[subsection]{Theorem}
\newtheorem{proposition}[subsection]{Proposition}
\newtheorem{lemma}[subsection]{Lemma}

\theoremstyle{definition}
\newtheorem{definition}[subsection]{Definition}
\newtheorem{example}[subsection]{Example}
\newtheorem{exercise}[subsection]{Exercise}
\newtheorem{situation}[subsection]{Situation}

\theoremstyle{remark}
\newtheorem{remark}[subsection]{Remark}
\newtheorem{remarks}[subsection]{Remarks}

\numberwithin{equation}{subsection}

% Macros
%
\def\lim{\mathop{\mathrm{lim}}\nolimits}
\def\colim{\mathop{\mathrm{colim}}\nolimits}
\def\Spec{\mathop{\mathrm{Spec}}}
\def\Hom{\mathop{\mathrm{Hom}}\nolimits}
\def\Ext{\mathop{\mathrm{Ext}}\nolimits}
\def\SheafHom{\mathop{\mathcal{H}\!\mathit{om}}\nolimits}
\def\SheafExt{\mathop{\mathcal{E}\!\mathit{xt}}\nolimits}
\def\Sch{\mathit{Sch}}
\def\Mor{\mathop{\mathrm{Mor}}\nolimits}
\def\Ob{\mathop{\mathrm{Ob}}\nolimits}
\def\Sh{\mathop{\mathit{Sh}}\nolimits}
\def\NL{\mathop{N\!L}\nolimits}
\def\CH{\mathop{\mathrm{CH}}\nolimits}
\def\proetale{{pro\text{-}\acute{e}tale}}
\def\etale{{\acute{e}tale}}
\def\QCoh{\mathit{QCoh}}
\def\Ker{\mathop{\mathrm{Ker}}}
\def\Im{\mathop{\mathrm{Im}}}
\def\Coker{\mathop{\mathrm{Coker}}}
\def\Coim{\mathop{\mathrm{Coim}}}

% Boxtimes
%
\DeclareMathSymbol{\boxtimes}{\mathbin}{AMSa}{"02}

%
% Macros for moduli stacks/spaces
%
\def\QCohstack{\mathcal{QC}\!\mathit{oh}}
\def\Cohstack{\mathcal{C}\!\mathit{oh}}
\def\Spacesstack{\mathcal{S}\!\mathit{paces}}
\def\Quotfunctor{\mathrm{Quot}}
\def\Hilbfunctor{\mathrm{Hilb}}
\def\Curvesstack{\mathcal{C}\!\mathit{urves}}
\def\Polarizedstack{\mathcal{P}\!\mathit{olarized}}
\def\Complexesstack{\mathcal{C}\!\mathit{omplexes}}
% \Pic is the operator that assigns to X its picard group, usage \Pic(X)
% \Picardstack_{X/B} denotes the Picard stack of X over B
% \Picardfunctor_{X/B} denotes the Picard functor of X over B
\def\Pic{\mathop{\mathrm{Pic}}\nolimits}
\def\Picardstack{\mathcal{P}\!\mathit{ic}}
\def\Picardfunctor{\mathrm{Pic}}
\def\Deformationcategory{\mathcal{D}\!\mathit{ef}}
\usepackage{cleveref}
\crefname{originalsection}{original section}{original sections}
\Crefname{originalsection}{Original section}{Original sections}
\makeatletter
\newlabel{duality-example-equidimensional-over-field}{{0AWJ}{source}{}{sourcepointer-1}{}}
\newlabel{duality-example-equidimensional-over-field@cref}{{[originalsection][0][]0AWJ}{[1][0][]source}{}{}{}}
\makeatother
\crefname{originalsection}{original section}{original sections}
\Crefname{originalsection}{Original section}{Original sections}
\makeatletter
\newlabel{duality-example-proper-over-local}{{0AWI}{source}{}{sourcepointer-2}{}}
\newlabel{duality-example-proper-over-local@cref}{{[originalsection][0][]0AWI}{[1][0][]source}{}{}{}}
\makeatother
\crefname{originallemma}{original lemma}{original lemmas}
\Crefname{originallemma}{Original lemma}{Original lemmas}
\makeatletter
\newlabel{duality-lemma-dualizing-module-proper-over-A}{{0AWP}{source}{}{sourcepointer-3}{}}
\newlabel{duality-lemma-dualizing-module-proper-over-A@cref}{{[originallemma][0][]0AWP}{[1][0][]source}{}{}{}}
\makeatother
\crefname{originallemma}{original lemma}{original lemmas}
\Crefname{originallemma}{Original lemma}{Original lemmas}
\makeatletter
\newlabel{duality-lemma-good-dualizing-normalized}{{0AWL}{source}{}{sourcepointer-4}{}}
\newlabel{duality-lemma-good-dualizing-normalized@cref}{{[originallemma][0][]0AWL}{[1][0][]source}{}{}{}}
\makeatother
\crefname{originallemma}{original lemma}{original lemmas}
\Crefname{originallemma}{Original lemma}{Original lemmas}
\makeatletter
\newlabel{duality-lemma-smooth-proper}{{0BRT}{source}{}{sourcepointer-5}{}}
\newlabel{duality-lemma-smooth-proper@cref}{{[originallemma][0][]0BRT}{[1][0][]source}{}{}{}}
\makeatother
\crefname{originallemma}{original lemma}{original lemmas}
\Crefname{originallemma}{Original lemma}{Original lemmas}
\makeatletter
\newlabel{duality-lemma-vanishing-good-dualizing}{{0AWN}{source}{}{sourcepointer-6}{}}
\newlabel{duality-lemma-vanishing-good-dualizing@cref}{{[originallemma][0][]0AWN}{[1][0][]source}{}{}{}}
\makeatother
\crefname{originalremark}{original remark}{original remarks}
\Crefname{originalremark}{Original remark}{Original remarks}
\makeatletter
\newlabel{duality-remark-relative-dualizing-complex}{{0B6S}{source}{}{sourcepointer-7}{}}
\newlabel{duality-remark-relative-dualizing-complex@cref}{{[originalremark][0][]0B6S}{[1][0][]source}{}{}{}}
\makeatother
\crefname{originalsection}{original section}{original sections}
\Crefname{originalsection}{Original section}{Original sections}
\makeatletter
\newlabel{duality-section-upper-shriek}{{0A9Y}{source}{}{sourcepointer-8}{}}
\newlabel{duality-section-upper-shriek@cref}{{[originalsection][0][]0A9Y}{[1][0][]source}{}{}{}}
\makeatother
\crefname{originallemma}{original lemma}{original lemmas}
\Crefname{originallemma}{Original lemma}{Original lemmas}
\makeatletter
\newlabel{dualizing-lemma-apply-CM}{{0AWS}{source}{}{sourcepointer-9}{}}
\newlabel{dualizing-lemma-apply-CM@cref}{{[originallemma][0][]0AWS}{[1][0][]source}{}{}{}}
\makeatother
\title{Existence and six functorial properties of the dualizing complex of a proper curve}
\author{The Stacks Project Authors}
\date{stacks-project:curves.tex:local-snapshot-20261012; GNU FDL 1.2 or later}
\begin{document}
\maketitle
\noindent\textit{Editorial note.} All delivered mathematical text is the author's own, copied from the pinned Stacks Project chapter file curves.tex; no statement or proof step was rewritten, reordered or removed.\par\medskip
\noindent\textit{Editorial note.} The counted proof is the author's complete argument: the relative-dualizing construction with adjointness, the vanishing and CM-support steps, the normalized local criterion, the unwound H\textasciicircum{}1-duality, and the smooth-case identification are all retained verbatim.\par\medskip
\noindent\textit{Editorial note.} The label lemma-duality-dim-1 names the counted statement, and the proof environment immediately following it in the source is delivered verbatim as its proof.\par\medskip
\noindent\textit{Editorial note.} The counted statement is carried by the author's own lemma environment in the chapter source.\par\medskip
\noindent\textit{Allowed standard prerequisites (editorial).}\begin{itemize}
\item The relative dualizing complex and its right adjoint. Duality for Schemes, Remark duality-remark-relative-dualizing-complex, Lemma duality-lemma-twisted-inverse-image, Section duality-section-upper-shriek. Defines omega\_X\textasciicircum{}bullet as a(O\_Spec(k)) for a the right adjoint of f\_*, gives property (4) by construction, and identifies a(O) with f\textasciicircum{}!(O) for proper f.
\item The proper and equidimensional framework examples. Duality for Schemes, Examples duality-example-proper-over-local and duality-example-equidimensional-over-field. Place omega\_X\textasciicircum{}bullet and omega\_X in the standard duality framework for a proper morphism to the spectrum of a field.
\item Vanishing, CM support, and local structure. Duality for Schemes, Lemmas duality-lemma-vanishing-good-dualizing and duality-lemma-good-dualizing-normalized; Dualizing Complexes, Lemma dualizing-lemma-apply-CM. Yield properties (1), (2) and the cohomologically graded local criterion of (3) through the normalized dualizing complex at closed points.
\item The dualizing module and the smooth case. Duality for Schemes, Lemmas duality-lemma-dualizing-module-proper-over-A and duality-lemma-smooth-proper. Give the H\textasciicircum{}1-duality of property (5) and the identification omega\_X = Omega\_\{X/k\} for smooth relative curves in property (6).
\end{itemize}\medskip
\ifdefined\hypertarget\hypertarget{sourcepointer-1}{}\fi
\noindent\textit{Source reference.} Original section 0AWJ: sources/stacks-chow/tags.txt, source lines 13722--13722. Prerequisite pointer invoked at proof line(s) 847; it names a standard fact of a sibling chapter that the counted argument uses, and every step of the counted proof built on it is retained verbatim in the delivered body.\par\medskip
\ifdefined\hypertarget\hypertarget{sourcepointer-2}{}\fi
\noindent\textit{Source reference.} Original section 0AWI: sources/stacks-chow/tags.txt, source lines 13721--13721. Prerequisite pointer invoked at proof line(s) 845; it names a standard fact of a sibling chapter that the counted argument uses, and every step of the counted proof built on it is retained verbatim in the delivered body.\par\medskip
\ifdefined\hypertarget\hypertarget{sourcepointer-3}{}\fi
\noindent\textit{Source reference.} Original lemma 0AWP: sources/stacks-chow/tags.txt, source lines 13727--13727. Prerequisite pointer invoked at proof line(s) 856; it names a standard fact of a sibling chapter that the counted argument uses, and every step of the counted proof built on it is retained verbatim in the delivered body.\par\medskip
\ifdefined\hypertarget\hypertarget{sourcepointer-4}{}\fi
\noindent\textit{Source reference.} Original lemma 0AWL: sources/stacks-chow/tags.txt, source lines 13724--13724. Prerequisite pointer invoked at proof line(s) 852; it names a standard fact of a sibling chapter that the counted argument uses, and every step of the counted proof built on it is retained verbatim in the delivered body.\par\medskip
\ifdefined\hypertarget\hypertarget{sourcepointer-5}{}\fi
\noindent\textit{Source reference.} Original lemma 0BRT: sources/stacks-chow/tags.txt, source lines 14781--14781. Prerequisite pointer invoked at proof line(s) 860; it names a standard fact of a sibling chapter that the counted argument uses, and every step of the counted proof built on it is retained verbatim in the delivered body.\par\medskip
\ifdefined\hypertarget\hypertarget{sourcepointer-6}{}\fi
\noindent\textit{Source reference.} Original lemma 0AWN: sources/stacks-chow/tags.txt, source lines 13726--13726. Prerequisite pointer invoked at proof line(s) 849; it names a standard fact of a sibling chapter that the counted argument uses, and every step of the counted proof built on it is retained verbatim in the delivered body.\par\medskip
\ifdefined\hypertarget\hypertarget{sourcepointer-7}{}\fi
\noindent\textit{Source reference.} Original remark 0B6S: sources/stacks-chow/tags.txt, source lines 14080--14080. Prerequisite pointer invoked at proof line(s) 837; it names a standard fact of a sibling chapter that the counted argument uses, and every step of the counted proof built on it is retained verbatim in the delivered body.\par\medskip
\ifdefined\hypertarget\hypertarget{sourcepointer-8}{}\fi
\noindent\textit{Source reference.} Original section 0A9Y: sources/stacks-chow/tags.txt, source lines 12966--12966. Prerequisite pointer invoked at proof line(s) 843; it names a standard fact of a sibling chapter that the counted argument uses, and every step of the counted proof built on it is retained verbatim in the delivered body.\par\medskip
\ifdefined\hypertarget\hypertarget{sourcepointer-9}{}\fi
\noindent\textit{Source reference.} Original lemma 0AWS: sources/stacks-chow/tags.txt, source lines 13730--13730. Prerequisite pointer invoked at proof line(s) 854; it names a standard fact of a sibling chapter that the counted argument uses, and every step of the counted proof built on it is retained verbatim in the delivered body.\par\medskip
\begin{lemma}
\label{lemma-duality-dim-1}
Let $X$ be a proper scheme of dimension $\leq 1$ over a field $k$.
There exists a dualizing complex $\omega_X^\bullet$ with the
following properties
\begin{enumerate}
\item $H^i(\omega_X^\bullet)$ is nonzero only for $i = -1, 0$,
\item $\omega_X = H^{-1}(\omega_X^\bullet)$
is a coherent Cohen-Macaulay module whose support is the
irreducible components of dimension $1$,
\item for $x \in X$ closed, the module $H^0(\omega_{X, x}^\bullet)$
is nonzero if and only if either
\begin{enumerate}
\item $\dim(\mathcal{O}_{X, x}) = 0$ or
\item $\dim(\mathcal{O}_{X, x}) = 1$
and $\mathcal{O}_{X, x}$ is not Cohen-Macaulay,
\end{enumerate}
\item for $K \in D_\QCoh(\mathcal{O}_X)$ there are functorial
isomorphisms\footnote{This property
characterizes $\omega_X^\bullet$ in $D_\QCoh(\mathcal{O}_X)$
up to unique isomorphism by the Yoneda lemma. Since $\omega_X^\bullet$
is in $D^b_{\textit{Coh}}(\mathcal{O}_X)$ in fact it suffices to consider
$K \in D^b_{\textit{Coh}}(\mathcal{O}_X)$.}
$$
\Ext^i_X(K, \omega_X^\bullet) = \Hom_k(H^{-i}(X, K), k)
$$
compatible with shifts and distinguished triangles,
\item there are functorial isomorphisms
$\Hom(\mathcal{F}, \omega_X) = \Hom_k(H^1(X, \mathcal{F}), k)$
for $\mathcal{F}$ quasi-coherent on $X$,
\item if $X \to \Spec(k)$ is smooth of relative dimension $1$,
then $\omega_X \cong \Omega_{X/k}$.
\end{enumerate}
\end{lemma}

\begin{proof}
Denote $f : X \to \Spec(k)$ the structure morphism.
We start with the relative dualizing complex
$$
\omega_X^\bullet = \omega_{X/k}^\bullet = a(\mathcal{O}_{\Spec(k)})
$$
as described in Duality for Schemes,
Remark \ref{duality-remark-relative-dualizing-complex}.
Then property (4) holds by construction as $a$ is the right
adjoint for $f_* : D_\QCoh(\mathcal{O}_X) \to D(\mathcal{O}_{\Spec(k)})$.
Since $f$ is proper we have
$f^!(\mathcal{O}_{\Spec(k)}) = a(\mathcal{O}_{\Spec(k)})$ by
definition, see
Duality for Schemes, Section \ref{duality-section-upper-shriek}.
Hence $\omega_X^\bullet$ and $\omega_X$ are as in
Duality for Schemes, Example \ref{duality-example-proper-over-local}
and as in
Duality for Schemes, Example \ref{duality-example-equidimensional-over-field}.
Parts (1) and (2) follow from
Duality for Schemes, Lemma \ref{duality-lemma-vanishing-good-dualizing}.
For a closed point $x \in X$ we see that $\omega_{X, x}^\bullet$ is a
normalized dualizing complex over $\mathcal{O}_{X, x}$, see
Duality for Schemes, Lemma \ref{duality-lemma-good-dualizing-normalized}.
Assertion (3) then follows from
Dualizing Complexes, Lemma \ref{dualizing-lemma-apply-CM}.
Assertion (5) follows from
Duality for Schemes, Lemma \ref{duality-lemma-dualizing-module-proper-over-A}
for coherent $\mathcal{F}$ and in general by unwinding
(4) for $K = \mathcal{F}[0]$ and $i = -1$.
Assertion (6) follows from Duality for Schemes,
Lemma \ref{duality-lemma-smooth-proper}.
\end{proof}

\end{document}
